\section{Introduction}\label{sec:introduction}

Termination of flips rules out an endless sequence of small birational improvements in the minimal model program. Since flips contract no divisors, their number cannot be bounded by counting divisors that disappear. In dimension four, exceptional surfaces bring places centered on surfaces and curves into the argument. On compact K\"ahler spaces, the contractions may be projective even when the ambient fourfold has no ample line bundle.

Shokurov's difficulty counts exceptional divisors with small discrepancy to measure the improvement of singularities \cite[Definition~2.15 and Corollaries~2.16--2.17]{ShokurovNonvanishing}. Kawamata, Matsuda, and Matsuki combined it with the rank of surface cycle classes to prove termination of terminal fourfold flips in the algebraic setting \cite[Theorem~5-1-15]{KawamataMatsudaMatsuki1987}. For boundaries, Fujino's corrected argument for canonical fourfold pairs exhibits the recurring contributions from successive blowups over a codimension-two center \cite[Example~2.1 and Theorems~5.1--5.2]{Fujino2005Addendum}. Alexeev, Hacon, and Kawamata developed a weighted difficulty that subtracts these contributions and tracks cycles on normalized boundary components \cite[Section~2]{AlexeevHaconKawamata2007}. Discrepancy improvement and the behavior of boundary cycles also guide the proof here.

Generalized pairs place part of the adjoint data on a higher birational model: a nef divisor there contributes its trace on the space carrying the boundary \cite[Definitions~1.4 and 4.1]{BirkarZhang2016}. For algebraic generalized pairs, Hacon and Moraga gave a termination reduction using weak Zariski decompositions \cite[Theorem~1]{HaconMoraga2020}. Termination for pseudo-effective generalized log canonical fourfolds with NQC nef data was proved by Chen and Tsakanikas \cite[Theorem~1.1]{ChenTsakanikas2023} and by Moraga \cite[Theorem~4.4]{Moraga2025}. Here NQC means a nonnegative real combination of nef \(\Q\)-Cartier divisors on the higher model. Han, Liu, and Zhuang proved termination of ordinary klt fourfold minimal model programs when the boundary is big over the fixed algebraic base \cite[Corollary~3.2(3)]{HanLiuZhuang2025}.

Hacon and Xie formulate a termination conjecture for generalized klt pairs with higher \((1,1)\)-class data on compact K\"ahler spaces \cite[Conjecture~1.15]{HaconXie2026}. The generalized klt case of the theorem below treats a four-dimensional rational-divisor form of that question under the stated projectivity and global Weil \(\Q\)-factoriality assumptions; the theorem also permits generalized log canonical singularities. Its ambient spaces may be nonprojective, and the prescribed small diagrams may vary their bases and choices under the hypotheses below.

A normal compact complex space \(X\) is \emph{globally Weil \(\Q\)-factorial} if every prime Weil divisor defined on all of \(X\) is \(\Q\)-Cartier and some positive reflexive power of its canonical sheaf is invertible. Projective morphisms of complex spaces are understood in the analytic sense: they admit a relatively ample holomorphic line bundle. A \(\Q\)-Cartier divisor is relatively ample if a positive Cartier multiple is relatively ample.

For a normal compact K\"ahler space \(X\), a rational nef b-divisor \(\M\) in this paper is represented by a \(\Q\)-Cartier divisor \(M'\) on a normal compact K\"ahler space \(X'\) with a projective bimeromorphic map \(\pi:X'\to X\). We require \(c_1(M')\) to lie in the closure of the K\"ahler cone in real Bott--Chern cohomology. On a model dominating \(X'\), the trace \(M_W\) is the pullback of \(M'\); on another model it is the pushforward from a common higher model. The divisor \(M'\) need not be effective.

We use compatible actual canonical divisor representatives on the birational models under consideration. Given an effective rational Weil divisor \(B\) of finite support for which \(D_X=K_X+B+M_X\) is \(\Q\)-Cartier, choose a smooth proper bimeromorphic model \(p:W\to X\) dominating \(X'\) and write
\begin{equation}\label{eq:discrepancy-definition}
 K_W+\Delta_W+M_W=p^*(K_X+B+M_X).
\end{equation}
A \emph{global divisorial place} is a prime divisor on a normal proper bimeromorphic model, with primes identified by strict transform on common higher models. Its log discrepancy is
\[
 a(E;X,B+\M)=1-\operatorname{coeff}_E(\Delta_W)
\]
on a model carrying \(E\). The pair is \emph{generalized log canonical} if this number is nonnegative for every global divisorial place, including the primes on \(X\). These places are defined on birational models; the definition does not depend on a field of global meromorphic functions.

\begin{theorem}\label{thm:main}
Let \(X_0\) be a normal irreducible globally Weil \(\Q\)-factorial compact K\"ahler fourfold. Let \(B_0\) be an effective rational Weil divisor of finite support with coefficients in \([0,1]\). Let \(\M\) be represented by an analytically nef \(\Q\)-Cartier divisor on a projective bimeromorphic normal compact K\"ahler model of \(X_0\). Fix compatible actual canonical divisor representatives on all birational models under consideration. Suppose that \(D_0=K_{X_0}+B_0+M_{X_0}\) is \(\Q\)-Cartier and that \((X_0,B_0+\M)\) is generalized log canonical.

Consider a sequence of normal irreducible globally Weil \(\Q\)-factorial compact K\"ahler fourfolds
\[
 X_0\dashrightarrow X_1\dashrightarrow X_2\dashrightarrow\cdots.
\]
On \(X_i\), let \(B_i\) be the strict transform of \(B_0\), let \(M_i=M_{X_i}\) be the trace of the same b-divisor, and assume that \(D_i=K_{X_i}+B_i+M_i\) is \(\Q\)-Cartier. Transport the canonical representatives under the induced identifications of prime divisors. Suppose that every step is given by a diagram
\begin{equation}\label{eq:flip-diagram}
\begin{tikzcd}[column sep=large]
 X_i \arrow[r,"f_i"] & Z_i & X_{i+1}\arrow[l,"f_i^+"']
\end{tikzcd}
\end{equation}
where \(Z_i\) is normal compact K\"ahler, both maps are projective small bimeromorphic morphisms with connected fibers, neither is an isomorphism, and
\[
 -D_i\text{ is }f_i\text{-ample},\qquad
 D_{i+1}\text{ is }f_i^+\text{-ample}.
\]
Here \emph{small} means that the exceptional locus contains no prime divisor. Then the sequence is finite.
\end{theorem}

We call a diagram satisfying the hypotheses of Theorem~\ref{thm:main} a \emph{flip} for the indicated adjoint. The theorem applies to arbitrary choices of such diagrams. Taking \(\M=0\) gives the ordinary log canonical case.

The proof uses the relative constructions and special termination for generalized pairs in \cite{KahlerAuxiliary}, whose precise forms we recall in Section~\ref{sec:preliminaries}.

\subsection*{Proof strategy}

We first prove termination for generalized klt pairs on models satisfying the stronger factoriality condition specified in Section~\ref{sec:preliminaries}. Discrepancy monotonicity and finiteness for the initial pair make the set of exceptional places of discrepancy at most one constant on a tail. Extracting it gives crepant models \(h_i:Y_i\to X_i\) whose remaining exceptional discrepancies are greater than one; we call them \emph{terminal junctions}. Each downstairs flip connects two junctions by a decrease of the extracted boundary coefficients followed by a finite, possibly empty, program of small steps. This construction adapts the projective argument of \cite[Section~3, Proposition~3.4]{ProjectiveLCFourfolds}.

The extracted coefficients may still vary. When an extracted divisor maps onto a surface, a transverse smooth surface cuts its general fiber into exceptional curves. Adjunction and negative definiteness bound the entries of the grouped intersection matrix. The fixed rational boundary and nef data then confine the discrepancies of extracted places with surface centers to a finite set. A downstairs flip strictly raises the discrepancy of a place whose center lies in its exceptional locus. It follows that, on every junction of a tail, each coefficient that still varies belongs to a divisor mapping onto a curve or a point of \(X_i\).

The branch bound of Theorem~\ref{thm:branch} applies to a reduced union \(Q\) of divisors on a terminal fourfold \(Y\) when a projective bimeromorphic morphism \(Y\to X\) of normal irreducible compact K\"ahler fourfolds sends \(Q\) to a set of dimension at most one. A surface in \(Q\) may have several prime divisors above it in the disjoint normalization of its components. The theorem bounds the total excess of branches by the ranks of surface cycle classes on those normalizations and the number of global ordinary discrepancy-two places centered on singular curves in \(Q\). It is formulated independently of any boundary or crepant equation.

The word \emph{global} is essential. Projective topology over the curve image relates the branch excess to local systems along the singular curves. Saito's projective analytic decomposition theorem and a Hodge-filtration argument on a compact K\"ahler resolution give the needed purity. A Quot parameter argument gives finite monodromy. On a Stein neighborhood of a compact core carrying all loops of the curve, we apply Fujino's local small \(\Q\)-factorialization theorem; fiber degrees and blowups then produce local ordinary discrepancy-two places. A canonical discrepancy comparison shows that they occur as divisors on one global resolution, and the retained loops keep those global divisors distinct. The normalization and parameter methods follow \cite[Sections~5--6]{ProjectiveLCFourfolds}; the analytic purity argument and the passage to distinct global divisors are given here.

Section~\ref{sec:difficulty} constructs the difficulty by counting exceptional places of discrepancy below two with ceiling weights, subtracting generic blowup contributions locally before summation, and adding cycle-rank corrections on normalized boundary components. The branch bound makes it nonnegative at junctions. Across a small step, the cycle-rank change cancels the change in the local subtractions, leaving a sum of nonnegative weight losses. This adapts the locally corrected difficulty in \cite[Section~4]{ProjectiveLCFourfolds}.

A surface in the positive exceptional locus of a downstairs flip forces a drop of at least one in the difficulty between its terminal junctions. Only finitely many flips can do this. Every remaining flip has an exceptional surface only on the negative side, so the rank of analytic surface classes on the ambient \(X_i\) decreases. This proves termination for generalized klt pairs on these auxiliary models. Special termination and deletion of the whole boundary floor reduce elementary generalized dlt programs to that case. We then use relative continuation and this gdlt termination to lift each diagram of Theorem~\ref{thm:main} to a finite program whose endpoint maps crepantly to the next downstairs model. A negative curve downstairs has a negative multisection upstairs, so each lift is nonempty. Their concatenation proves Theorem~\ref{thm:main}.
